% `tex.web` §1035's `ins_disc`, and why a compound word may break after its
% hyphen without anyone writing `\-`:
%
%     if ins_disc then
%       begin ins_disc:=false;
%       if mode>0 then tail_append(new_disc);
%       end;
%
% The character just appended is the current font's \hyphenchar, so an empty
% discretionary goes after it. We had none of it: every hyphenated compound in
% every document was one unbreakable word, which forces a worse line break
% wherever one falls near it.
%
% Three conditions, and each is a separate way to get it wrong.
%
%   * It is the *font's* hyphen character, not the literal `-`. Set
%     \hyphenchar to `Q` and it is `Q` that gets the break and `-` that does
%     not — which is how a font whose hyphen sits elsewhere still works.
%
%   * A negative \hyphenchar gets none at all. That is the same test §896
%     applies before hyphenating a word, and setting \hyphenchar to -1 is how
%     a document turns both off together.
%
%   * `mode>0`: a real paragraph. An \hbox has no lines to break, so a hyphen
%     inside one gets nothing — including an \hbox sitting in a paragraph,
%     where the hyphens on either side of it are treated differently.
\catcode`\{=1 \catcode`\}=2
\tracingonline=1 \showboxbreadth=99 \showboxdepth=99
\hsize=200pt \parindent=0pt \vsize=200pt
\font\r=cmr10 \r

% The ordinary case.
\hyphenchar\r=`\-
\setbox0\vbox{A-B\par}
\showbox0

% The font's hyphen is `Q`, so `-` is an ordinary character and `Q` is not.
\hyphenchar\r=`\Q
\setbox1\vbox{A-BQC\par}
\showbox1

% Off.
\hyphenchar\r=-1
\setbox2\vbox{A-BQC\par}
\showbox2

% `\char` appends a character like any other, and gets the same break.
\hyphenchar\r=`\-
\setbox3\vbox{A\char`\- B\par}
\showbox3

% An \hbox gets none, in a paragraph or out of one.
\setbox4\hbox{A-B}
\showbox4
\setbox5\vbox{\hbox{A-B}-C\par}
\showbox5

% ...and `\-` still contributes its own, which is not this one: it carries the
% \hyphenchar in its pre-break text, and the character after it is not a
% hyphen, so nothing follows it.
\setbox6\vbox{A\-B\par}
\showbox6
\end
